\chapter{Cibles, étiquettes et poids d'échantillon}\label{ch:ml:targets-labels-and-sample-weights}

Un classifieur est entraîné pour dire si un actif sera en hausse dans dix jours, et il fait un peu mieux qu'un pile ou
face. Le trader qui va l'utiliser demande ce qu'il fait du stop-loss~: chaque position est clôturée si elle perd une volatilité de dix jours,
et prise de profit si elle gagne autant. Sur vingt actifs synthétiques observés pendant dix ans, une transaction sur cinq atteint
son stop avant le dixième jour~; l'étiquette apprise par le classifieur n'a jamais tenu compte du stop. Et l'\omterm{def:ml:why-financial-machine-learning-is-different:sets}{ensemble d'entraînement}
qui semble compter 49\,000 exemples ne contient, une fois le chevauchement des fenêtres de dix jours comptabilisé, que l'information de 4\,900. Ce chapitre construit
des étiquettes qui décrivent la transaction, mesure le degré de chevauchement des étiquettes, et transforme ce chevauchement en \omterm{def:ml:targets-labels-and-sample-weights:weights}{poids d'échantillon}
et en statistiques honnêtes.

\section{Étiquettes à horizon fixe et leurs défauts}

Les données du chapitre proviennent de \texttt{firm.mlsynth.series}~: vingt actifs, 2\,520 jours, une volatilité quotidienne de 1.5~\% avec
un regroupement de type GARCH et des chocs de Student-$t$, ainsi qu'une dérive planifiée, persistante, de demi-vie 20 jours, dont la corrélation
avec le rendement du jour suivant vaut 0.04. Les variables, connues à chaque clôture, sont une lecture bruitée de la dérive, les rendements passés sur 5 et
20 jours, un rapport entre la volatilité réalisée courte et la volatilité réalisée longue, et trois colonnes de pur bruit. Un événement est chaque jour
de chaque actif~; la transaction principale est longue lorsque la lecture de la dérive est positive, et courte dans le cas contraire.

\begin{definition}[Étiquette à horizon fixe]\label{def:ml:targets-labels-and-sample-weights:fixed}
Soit un événement à la barre $t_0$ et un horizon $h$. L'\emph{étiquette à horizon fixe}\index{etiquette@étiquette à horizon fixe} (fixed-horizon label) vaut
$\sign(\sum_{s=t_0+1}^{t_0+h}R_s)$, ou le rendement lui-même pour la régression, réglée en $t_1 = t_0 + h$.
\end{definition}

L'\omterm{def:ml:targets-labels-and-sample-weights:fixed}{étiquette à horizon fixe} est le rendement à terme du livre~7 (chapitre~6) transformé en classe. Elle présente trois défauts. Elle ignore le chemin~: une position
sortie par le stop au troisième jour et une position qui a finalement atteint le dixième jour reçoivent la même étiquette si le cours est revenu. Elle ignore la volatilité~: un mouvement de 1~\% est du bruit un jour volatil et une information
un jour calme. Et elle se termine à un instant fixe, si bien que les étiquettes d'événements consécutifs se chevauchent sur $h - 1$ de leurs $h$ barres. Les deux premiers défauts
sont corrigés par l'\omterm{def:ml:targets-labels-and-sample-weights:triple}{étiquette à triple barrière} ci-dessous~; le troisième
par le comptage (la \cref{sec:ml:labels:overlap}).

\section{Étiquettes à barrières et méta-étiquetage}

\begin{definition}[Étiquette à triple barrière]\label{def:ml:targets-labels-and-sample-weights:triple}
Pour un événement en $t_0$ de côté de position $\pm1$, des largeurs $w^+, w^- > 0$ et un horizon $h$, l'\emph{étiquette à triple
barrière}\index{etiquette@étiquette à triple barrière} (triple-barrier label) se règle à la première barre $t_1$ où le rendement logarithmique cumulé de la position atteint
$+w^+$ (prise de profit), $-w^-$ (stop-loss), ou $t_1 = t_0 + h$ (la barrière verticale)~; elle est le signe du rendement de la position
en $t_1$. Les largeurs sont en général mises à l'échelle par la volatilité connue en $t_0$~: $w = k\hat\sigma_{t_0}\sqrt h$
(López de Prado, 2018).
\end{definition}

Avec $k = 1$ et $h = 10$, la transaction longue ou courte du chapitre atteint sa barrière de profit en premier dans 22.7~\% des cas,
son stop dans 20.3~\% des cas, et la barrière verticale sinon~; une étiquette dure en moyenne 8.3 jours au lieu de 10.
La \Cref{fig:ml:labels:path} montre un événement du premier actif~: une position longue sortie au troisième jour à $-5.04~\%$ (la
largeur était de 4.87~\%), alors que le cours a fini le dixième jour 1.84~\% plus haut. L'\omterm{def:ml:targets-labels-and-sample-weights:fixed}{étiquette à horizon fixe} déclare cette transaction gagnante.
C'est un cas parmi beaucoup~: parmi les transactions sorties par le stop, 5.3~\% se terminent positivement au dixième jour. Les signes ne concordent pas pour 1.9~\% de l'ensemble des événements~;
ce qui diffère partout, c'est la taille, puisque les rendements de l'étiquette à barrières sont tronqués aux barrières
comme ceux de la transaction.

\begin{omfigure}
\begin{tikzpicture}
\begin{axis}[width=10.5cm,height=5.6cm,xlabel={jours après l'événement},ylabel={rendement cumulé (\%)},
  tick label style={font=\small},label style={font=\small},xmin=-0.3,xmax=10.3,ymin=-7,ymax=6.5,
  xtick={0,1,2,3,4,5,6,7,8,9,10},grid=major,grid style={gray!25},table/col sep=comma]
\addplot[omDef,very thick,mark=*,mark size=1.6pt] table[x=day,y=ret]{figdata/ml/02-targets-labels-and-sample-weights/path.csv};
\addplot[omProp,thick,dashed] table[x=day,y=up]{figdata/ml/02-targets-labels-and-sample-weights/path.csv};
\addplot[omThm,thick,dashed] table[x=day,y=down]{figdata/ml/02-targets-labels-and-sample-weights/path.csv};
\draw[black!70,thick] (axis cs:10,-7) -- (axis cs:10,6.5);
\node[font=\footnotesize,omProp,anchor=south west] at (axis cs:0.2,4.87) {prise de profit};
\node[font=\footnotesize,omThm,anchor=north] at (axis cs:8.2,-5.0) {stop-loss};
\node[font=\footnotesize,black!70,anchor=east] at (axis cs:9.9,5.6) {barrière verticale};
\draw[omThm,thick] (axis cs:3,-5.04) circle[radius=4pt];
\node[font=\footnotesize,omThm,anchor=west] at (axis cs:3.3,-6.2) {stop touché au jour 3};
\end{axis}
\end{tikzpicture}

\omcaption{Un événement de l'actif~0, une position longue~: les barrières valent une volatilité de dix jours, soit 4.87~\%, connue à l'événement.
Le stop est touché au troisième jour~; le rendement sur dix jours se termine à $+1.84\%$. Données~: \texttt{firm.mlsynth.series}, graine~1,
via \texttt{firm.labeling.triple\_barrier}.}\label{fig:ml:labels:path}
\end{omfigure}

Le côté d'une transaction et sa taille sont deux questions différentes. Un modèle primaire (une règle, le signal d'un chercheur, un choix discrétionnaire)
peut être bon pour le côté et médiocre pour savoir quand il a raison.

\begin{definition}[Méta-étiquetage]\label{def:ml:targets-labels-and-sample-weights:meta}
Le \emph{méta-étiquetage}\index{méta-étiquetage} (meta-labelling) étiquette chaque événement d'un modèle primaire avec 1 si la transaction dans le sens du modèle primaire a rapporté
(par exemple, si son rendement à triple barrière est positif) et avec 0 sinon, puis entraîne un modèle secondaire sur ces étiquettes
pour estimer la probabilité que l'appel primaire soit juste~; cette probabilité sert ensuite à filtrer ou à dimensionner la transaction
(López de Prado, 2018~; Joubert, 2022).
\end{definition}

Le \omterm{def:ml:targets-labels-and-sample-weights:meta}{méta-étiquetage} de la transaction du chapitre à l'aide d'arbres à gradient boosté (la \cref{fig:ml:labels:meta}), entraîné sur les premiers 60~\%
des dates et évalué sur les derniers 40~\%, fait passer le taux de réussite des transactions retenues de 53.0~\% à 54.5~\% et le rendement moyen
par transaction de 0.064 à 0.097 largeur de barrière lorsqu'on conserve le tiers des transactions dont le modèle est le plus sûr. La $t$-statistique ajustée du chevauchement
tombe de 3.48 à 3.19~: les transactions retenues sont meilleures, mais il y en a moins.
Le \omterm{def:ml:targets-labels-and-sample-weights:meta}{méta-étiquetage} réaffecte le risque vers les meilleures transactions~; il ne crée pas un signal qui manque au modèle primaire.

\begin{omfigure}
\begin{tikzpicture}
\begin{groupplot}[group style={group size=2 by 1,horizontal sep=1.7cm},width=6.6cm,height=5.2cm,
  xlabel={\small transactions conservées (\%)},tick label style={font=\footnotesize},label style={font=\small},x dir=reverse,
  xmin=5,xmax=105,grid=major,grid style={gray!25},table/col sep=comma]
\nextgroupplot[ylabel={rendement moyen par transaction (largeurs)},ymin=0.05,ymax=0.13,
  yticklabel style={/pgf/number format/fixed,/pgf/number format/precision=2},scaled y ticks=false]
\addplot[omDef,very thick,mark=*] table[x=share,y=mean]{figdata/ml/02-targets-labels-and-sample-weights/meta.csv};
\nextgroupplot[ylabel={$t$ ajustée du chevauchement},ymin=2,ymax=3.8]
\addplot[omThm,very thick,mark=square*] table[x=share,y=t]{figdata/ml/02-targets-labels-and-sample-weights/meta.csv};
\end{groupplot}
\end{tikzpicture}

\omcaption{\omterm{def:ml:targets-labels-and-sample-weights:meta}{Méta-étiquetage} de la transaction primaire~: on conserve les transactions dont la probabilité de réussite estimée dépasse 0.50
à 0.55 (100~\% = le modèle primaire seul). À gauche, le rendement moyen par transaction en largeurs de barrière~; à droite, la $t$-statistique
avec le nombre de transactions ajusté du chevauchement. Hors échantillon, derniers 40~\% des dates. Données~: \texttt{ml\_labels.meta}.}\label{fig:ml:labels:meta}
\end{omfigure}

\section{Chevauchement, concurrence et unicité}\label{sec:ml:labels:overlap}

\begin{definition}[Concurrence des étiquettes, unicité moyenne]\label{def:ml:targets-labels-and-sample-weights:uniqueness}
L'étiquette $i$ couvre les barres $(t_{0,i}, t_{1,i}]$. La \emph{concurrence des étiquettes}\index{concurrence des étiquettes} (label concurrency) $c_t$ de la barre $t$
est le nombre d'étiquettes dont l'intervalle la contient. L'\emph{unicité moyenne}\index{unicité moyenne} (average uniqueness) de l'étiquette $i$ est
$\bar u_i = \frac{1}{t_{1,i} - t_{0,i}}\sum_{t\in(t_{0,i}, t_{1,i}]}1/c_t$~: la part de son information qu'elle détient seule.
\end{definition}

La \Cref{fig:ml:labels:toy} est le plus petit exemple~: trois étiquettes sur sept barres, dont deux se chevauchent sur deux barres.
Le livre~7 a nommé le phénomène (chevauchement des étiquettes, chapitre~20) et en a retiré la fuite de la validation croisée par la purge~; ici
on le mesure, parce qu'il modifie la valeur de l'échantillon.

\begin{omfigure}
\begin{tikzpicture}[x=1.15cm,y=0.62cm,font=\small]
\foreach \b in {1,2,3,4,5,6,7} { \draw[gray!50] (\b-0.5,-0.2) -- (\b-0.5,4.4); \node at (\b,4.1) {\b}; }
\draw[gray!50] (7.5,-0.2) -- (7.5,4.4);
\node[anchor=east] at (0.3,4.1) {barre};
\draw[omDef,line width=4pt] (0.55,3) -- (3.45,3); \node[anchor=east] at (0.3,3) {étiquette A};
\draw[omProp,line width=4pt] (1.55,2) -- (4.45,2); \node[anchor=east] at (0.3,2) {étiquette B};
\draw[omThm,line width=4pt] (5.55,1) -- (7.45,1); \node[anchor=east] at (0.3,1) {étiquette C};
\node[anchor=east] at (0.3,0) {$c_t$};
\foreach \b/\c in {1/1,2/2,3/2,4/1,5/0,6/1,7/1} { \node at (\b,0) {\c}; }
\node[anchor=west,align=left,font=\footnotesize] at (7.8,2) {$\bar u_A = (1+\tfrac12+\tfrac12)/3 = 0.67$\\
$\bar u_B = (\tfrac12+\tfrac12+1)/3 = 0.67$\\ $\bar u_C = 1$};
\end{tikzpicture}

\omcaption{Concurrence et \omterm{def:ml:targets-labels-and-sample-weights:uniqueness}{unicité moyenne} sur sept barres~: les étiquettes A et B partagent les barres 2 et 3~; C est seule. Les trois
étiquettes valent $0.67 + 0.67 + 1 = 2.33$ étiquettes indépendantes, soit le nombre de barres qu'elles couvrent divisé par leur longueur
seulement lorsque toutes les longueurs sont égales (\cref{prop:ml:labels:sum}).}\label{fig:ml:labels:toy}
\end{omfigure}

\begin{proposition}[Valeur des étiquettes qui se chevauchent]\label{prop:ml:labels:sum}
\begin{enumerate}
\item Si chaque étiquette couvre $h$ barres, $\sum_i\bar u_i = N_{\mathrm{cov}}/h$, où $N_{\mathrm{cov}}$ est le nombre de
barres couvertes par au moins une étiquette. Avec un événement sur chaque barre, $\bar u_i = 1/h$ sauf aux extrémités.
\item Si les rendements sont non corrélés, de variance $\sigma^2$, et si les $n$ étiquettes sont les sommes chevauchantes de $h$ barres d'une même série, la
moyenne empirique des étiquettes a pour variance environ $h^2\sigma^2/n$, soit $h$ fois le $h\sigma^2/n$ que supposerait leur traitement
comme indépendantes~: une statistique $t$ calculée sur $n$ lignes est gonflée d'environ $\sqrt h$.
\end{enumerate}
\end{proposition}

\begin{proof}
(1) $\sum_i\bar u_i = \frac1h\sum_i\sum_{t\in i}1/c_t = \frac1h\sum_t\sum_{i\ni t}1/c_t = \frac1h\#\{t: c_t\ge1\}$,
puisque la somme intérieure comporte $c_t$ termes égaux à $1/c_t$. (2) Loin des extrémités, chaque rendement entre dans $h$ des $n$ sommes, si bien que
la moyenne des sommes vaut environ $\frac hn\sum_tR_t$, de variance $\frac{h^2}{n^2}\,n\sigma^2 = h^2\sigma^2/n$.
\end{proof}

Sur les données du chapitre, les \omterm{def:ml:targets-labels-and-sample-weights:fixed}{étiquettes à horizon fixe} ont une concurrence exactement égale à 10 et une \omterm{def:ml:targets-labels-and-sample-weights:uniqueness}{unicité moyenne} de 0.100~:
48\,980 étiquettes valent 4\,916, comme le prédit la partie 1 ($20\times2\,458/10$). Les étiquettes par barrières, plus courtes, ont une unicité
de 0.123 et valent 6\,040. La rentabilité moyenne de la transaction primaire a une statistique $t$ de 9.4 calculée sur les lignes et
3.3 lorsque le nombre de transactions est remplacé par la somme des unicités~: le rapport, 2.85, vaut $\sqrt{1/0.123}$.

\section{Poids d'échantillon et bootstrap séquentiel}

\begin{definition}[Poids d'échantillon, poids à décroissance temporelle]\label{def:ml:targets-labels-and-sample-weights:weights}
Un \emph{poids d'échantillon}\index{poids d'échantillon} (sample weight) $\omega_i\ge0$ pondère le terme de l'observation $i$ dans la perte empirique,
$\mathcal L_n(\theta) = \sum_i\omega_i\ell(y_i, f_\theta(x_i))/\sum_i\omega_i$. Trois sont courants sur les données de marché~:
les poids d'unicité $\omega_i = \bar u_i$~; les poids d'attribution du rendement, $\omega_i\propto|\sum_{t\in i}R_t/c_t|$, qui
favorisent les étiquettes dont ils possèdent les mouvements~; et un \emph{poids à décroissance temporelle}\index{poids à décroissance temporelle} (time-decay weight), qui décroît linéairement de 1
pour l'étiquette la plus récente à un plancher pour la plus ancienne, en unités d'unicité cumulée, de sorte que les anciennes données comptent moins.
\end{definition}

\begin{definition}[Bootstrap séquentiel]\label{def:ml:targets-labels-and-sample-weights:seqboot}
Le \emph{bootstrap séquentiel}\index{bootstrap séquentiel} (sequential bootstrap) tire les étiquettes une à une, avec remise, chacune avec
une probabilité proportionnelle à l'\omterm{def:ml:targets-labels-and-sample-weights:uniqueness}{unicité moyenne} qu'elle aurait compte tenu des étiquettes déjà tirées~; les étiquettes qui se chevauchent
deviennent moins probables dès qu'une d'entre elles figure dans l'échantillon.
\end{definition}

L'argument en faveur de ces outils tient à ceci~: un ensemble d'arbres agrégés dont les échantillons bootstrap regorgent de quasi-doublons produit des arbres corrélés,
chacun ajusté sur les mêmes quelques observations indépendantes. En tirant 40 des 400 premières étiquettes à dix jours d'un actif
(environ leur nombre effectif), les échantillons du bootstrap classique ont une \omterm{def:ml:targets-labels-and-sample-weights:uniqueness}{unicité moyenne} de 0.65 et ceux du \omterm{def:ml:targets-labels-and-sample-weights:seqboot}{bootstrap séquentiel}
de 0.69. Que cela importe pour la prévision est une question empirique, et sur ces données la réponse est non
(\cref{tab:ml:labels:bagging})~: des arbres agrégés entraînés sur des échantillons bootstrap classiques, sur des échantillons aussi petits que
l'\omterm{def:ml:targets-labels-and-sample-weights:uniqueness}{unicité moyenne} multipliée par la taille de l'échantillon, et sur ces échantillons pondérés par l'unicité, ont la même AUC hors échantillon dans trois
marchés simulés, à $\pm0.007$ les uns des autres et sans ordre constant. Les arbres étaient régularisés (au moins 50
étiquettes par feuille, 5 dans les petits échantillons), et aucune des variables ne permet d'identifier une date, ce qui rend
précisément possible la fuite des quasi-doublons. Là où le chevauchement mord à coup sûr, c'est dans toute statistique qui compte des lignes~:
erreurs types, statistiques $t$, intervalles de bootstrap et scores hors sac.

\begin{table}[htbp]
\centering\small
\begin{tabular}{@{}lcccc@{}}
\toprule
échantillons bootstrap & marché 1 & marché 2 & marché 3 & moyenne\\
\midrule
classique, taille pleine & 0.504 & 0.517 & 0.509 & 0.510\\
taille à \omterm{def:ml:targets-labels-and-sample-weights:uniqueness}{unicité moyenne} & 0.511 & 0.512 & 0.506 & 0.510\\
idem, pondéré par l'unicité & 0.506 & 0.511 & 0.510 & 0.509\\
\bottomrule
\end{tabular}
\caption{AUC hors échantillon de 30 arbres agrégés prédisant le signe de la rentabilité à dix jours, entraînés sur les premiers 60~\% des
dates (purgées) de trois marchés simulés. Données~: \texttt{ml\_labels.bagging\_table}.}\label{tab:ml:labels:bagging}
\end{table}

\begin{method}[Étiquettes et poids pour une nouvelle étude]\label{met:ml:labels:method}
\begin{enumerate}
\item Écrire d'abord la transaction (côté, entrée, sorties, horizon), puis l'étiquette qui la règle~: barrières au stop et à l'objectif
de la transaction, mises à l'échelle par la volatilité connue à l'événement.
\item Conserver chaque étiquette avec son $t_0$ et son $t_1$~; calculer la concurrence et l'unicité par actif.
\item Indiquer la taille effective de l'échantillon $\sum_i\bar u_i$ à côté du nombre de lignes, et calculer chaque erreur type avec elle
(ou avec un bootstrap par blocs du livre~4, chapitre~13).
\item Essayer les poids d'unicité et à décroissance temporelle comme \omterm{def:ml:why-financial-machine-learning-is-different:sets}{hyperparamètres}, sur des plis de validation purgés (chapitre~3)~; ne les garder
que s'ils aident là.
\end{enumerate}
\end{method}

\section{Tutoriel~: des étiquettes qui décrivent la transaction}\label{tut:ml:targets-labels-and-sample-weights:tut}

\begin{tutorial}
\textbf{Objectif.} Étiqueter les événements quotidiens de vingt actifs avec des \omterm{def:ml:targets-labels-and-sample-weights:fixed}{étiquettes à horizon fixe} et à triple barrière, mesurer le chevauchement, comparer
les schémas de bootstrap et méta-étiqueter la transaction primaire. \textbf{État final~:} les chiffres du chapitre,
les \cref{fig:ml:labels:path,fig:ml:labels:meta}, le \cref{tab:ml:labels:bagging}.
\begin{tutsteps}
\item \textbf{Étiquettes par barrières}~: le premier contact de l'une des barrières, ou la barrière verticale.
\omcode{code/firm/labeling/firm_labeling.py}{41}{70}{L'étiquette à triple barrière.}\label{lst:ml:labels:triple}
\item \textbf{Concurrence et unicité} par sommes cumulées.
\omcode{code/firm/labeling/firm_labeling.py}{77}{92}{Concurrence et unicité moyenne.}\label{lst:ml:labels:uniq}
\item \textbf{Le \omterm{def:ml:targets-labels-and-sample-weights:seqboot}{bootstrap séquentiel}}~: chaque tirage favorise les étiquettes qui se chevauchent le moins avec celles déjà tirées.
\omcode{code/firm/labeling/firm_labeling.py}{110}{123}{Le bootstrap séquentiel.}\label{lst:ml:labels:seq}
\item \textbf{Exécuter} \texttt{ml\_labels.label\_stats()}, \texttt{bagging\_table()}, \texttt{bootstrap\_uniqueness()},
\texttt{meta()} et \texttt{fig\_labels.py}.
\end{tutsteps}
\textbf{Pour aller plus loin.} Placer les barrières à deux volatilités et observer comment changent la part des sorties par barrière verticale
et l'unicité~; méta-étiqueter avec une régression logistique au lieu des arbres boostés.
\end{tutorial}

\section{Construction~: étiquettes et poids d'échantillon}\label{bld:ml:targets-labels-and-sample-weights:labeling}

\begin{build}
\bfield{Rôle} Chaque étude étiquette ses événements comme le desk les traite, et connaît le nombre
d'observations indépendantes dont elle dispose.
\bfield{Interface} \texttt{fixed\_horizon(r, t0, h)}, \texttt{triple\_barrier(r, t0, up, down, h, side)},
\texttt{vol\_widths(sigma, t0, h, k)}, \texttt{meta\_labels(side, ret)}, \texttt{concurrency(t0, t1, n)},
\texttt{average\_uniqueness(t0, t1, n)}, \texttt{attribution\_weights(t0, t1, r, n)}, \texttt{time\_decay(u, last)},
\texttt{sequential\_bootstrap(t0, t1, n, size, rng)}.
\bfield{Règles} Chaque étiquette porte $t_0$ et $t_1$~; la clôture de la barre $t_0$ est l'instant de décision et l'étiquette utilise les barres
$t_0 + 1$ à $t_1$~; les largeurs des barrières n'utilisent que l'information connue en $t_0$.
\bfield{Tests d'acceptation} \texttt{code/firm/labeling/tests/}~: étiquettes et touches vérifiées à la main, pour une position longue et une position courte~;
concurrence, unicité, poids d'attribution et de décroissance sur un exemple de trois étiquettes~; le \omterm{def:ml:targets-labels-and-sample-weights:seqboot}{bootstrap séquentiel} élève
l'\omterm{def:ml:targets-labels-and-sample-weights:uniqueness}{unicité moyenne} des petits échantillons.
\bfield{Pour aller plus loin} Étiquettes de balayage de tendance (l'horizon présentant la tendance la plus significative)~; étiquettes pour les exécutions du market making
(chapitre~29)~; une recherche vectorisée de barrières sur de nombreux événements à la fois.
\end{build}

\begin{omsources}
\item M.~López de Prado, \emph{Advances in Financial Machine Learning}, Wiley, 2018, chapitres 3--4.
\item J.~F.~Joubert, «~Meta-labeling: theory and framework~», \emph{Journal of Financial Data Science} 4(3), 2022.
\item L.~Breiman, «~Bagging predictors~», \emph{Machine Learning} 24, 1996.
\end{omsources}

\section{Exercices}

\begin{exercise}[$\star$]\label{exo:ml:targets-labels-and-sample-weights:1}
Des étiquettes de 5 barres débutent sur chaque barre d'une série de 1\,000 barres (les dernières étant tronquées à la fin). À peu près combien
d'étiquettes indépendantes représentent-elles, et quelle est leur \omterm{def:ml:targets-labels-and-sample-weights:uniqueness}{unicité moyenne} hors des extrémités~?
\end{exercise}

\begin{exercise}[$\star$]\label{exo:ml:targets-labels-and-sample-weights:2}
Quatre étiquettes couvrent les barres $(0,2]$, $(1,4]$, $(2,4]$ et $(6,8]$. Écrire la concurrence des barres 1 à 8 et
l'\omterm{def:ml:targets-labels-and-sample-weights:uniqueness}{unicité moyenne} de chaque étiquette.
\end{exercise}

\begin{exercise}[$\star$]\label{exo:ml:targets-labels-and-sample-weights:3}
La volatilité quotidienne connue à un événement est de 1.2~\%. Quelles sont les largeurs des barrières pour $k = 1.5$ et un horizon de 16 jours~?
\end{exercise}

\begin{exercise}[$\star\star$]\label{exo:ml:targets-labels-and-sample-weights:4}
La transaction moyenne d'une stratégie a une $t$-statistique de 6.1, calculée sur 25\,000 transactions de 20 jours qui se chevauchent, ouvertes chaque jour.
Quelle $t$-statistique faut-il rapporter, et pourquoi~?
\end{exercise}

\begin{exercise}[$\star\star$]\label{exo:ml:targets-labels-and-sample-weights:5}
Pourquoi le \omterm{def:ml:targets-labels-and-sample-weights:meta}{méta-étiquetage} augmente-t-il le rendement moyen par transaction sans augmenter la $t$-statistique de l'ensemble du portefeuille, dans
la \cref{fig:ml:labels:meta}~? Quand augmenterait-il les deux~?
\end{exercise}

\begin{exercise}[$\star\star$]\label{exo:ml:targets-labels-and-sample-weights:6}
\emph{Trouver la faille.} «~Nos largeurs de barrières sont fixées à un écart type des rendements quotidiens de chaque actif sur l'ensemble de
l'échantillon, multiplié par la racine carrée de l'horizon.~»
\end{exercise}

\begin{exercise}[$\star\star\star$]\label{exo:ml:targets-labels-and-sample-weights:7}
\emph{Programmation.} Reconstruire \texttt{ml\_labels.events} avec des barrières à deux volatilités sur dix jours ($k = 2$). Rapporter les parts
des sorties à profit, à stop et par barrière verticale, ainsi que l'\omterm{def:ml:targets-labels-and-sample-weights:uniqueness}{unicité moyenne} des étiquettes à barrières, et expliquer le changement.
\end{exercise}

\begin{exercise}[$\star\star\star$]\label{exo:ml:targets-labels-and-sample-weights:8}
Avec des poids de décroissance temporelle décroissant linéairement de 1 à 0.5 en unicité cumulée, et chaque étiquette d'unicité 0.1, quel est le poids de l'étiquette
à mi-parcours de l'échantillon~? Montrer que la moyenne des poids vaut 0.75 quelle que soit l'unicité, si celle-ci
est constante, et dire ce qui change si elle ne l'est pas.
\end{exercise}

\section{Problème~: Cinq cents jours qui se chevauchent}

\begin{problem}[{Problème du week-end --- ce que vaut un ensemble d'entraînement}]\label{pb:ml:targets-labels-and-sample-weights:1}
Les vingt actifs du chapitre, leurs événements quotidiens et la transaction primaire, longue ou courte.

\medskip\noindent\textbf{Partie I --- Deux étiquettes.}
\begin{enumerate}
\item Avec quelle fréquence la transaction atteint-elle sa barrière de profit, son stop, la barrière verticale~?
\item Quelle est la durée moyenne d'une étiquette à barrière, et pourquoi est-elle inférieure à dix jours~?
\item Dans la \cref{fig:ml:labels:path}, que disent l'\omterm{def:ml:targets-labels-and-sample-weights:fixed}{étiquette à horizon fixe} et l'étiquette à barrière, et laquelle décrit la transaction~?
\item Quelle part de l'ensemble des événements porte des étiquettes de signes opposés, et pourquoi cela ne constitue-t-il pas toute la différence~?
\end{enumerate}

\medskip\noindent\textbf{Partie II --- Chevauchement.}
\begin{enumerate}[resume]
\item Quelles sont la concurrence et l'\omterm{def:ml:targets-labels-and-sample-weights:uniqueness}{unicité moyenne} des \omterm{def:ml:targets-labels-and-sample-weights:fixed}{étiquettes à horizon fixe}~?
\item Combien d'étiquettes indépendantes valent les 48\,980, selon la \cref{prop:ml:labels:sum}~? Et les étiquettes à barrières~?
\item Donner la $t$-statistique de la transaction sur les lignes et celle ajustée du chevauchement, puis expliquer leur rapport.
\item Quelles statistiques d'une étude sont fausses si l'on ignore le chevauchement~?
\end{enumerate}

\medskip\noindent\textbf{Partie III --- Pondérations et bootstraps.}
\begin{enumerate}[resume]
\item Quelle \omterm{def:ml:targets-labels-and-sample-weights:uniqueness}{unicité moyenne} ont les échantillons de 40 étiquettes issus d'un bootstrap classique et d'un \omterm{def:ml:targets-labels-and-sample-weights:seqboot}{bootstrap séquentiel}~?
\item Quelle AUC atteignent les trois schémas de bagging, marché par marché~?
\item Pourquoi l'unicité n'a-t-elle rien changé à l'AUC ici~?
\item Concevoir un cadre où elle le ferait.
\end{enumerate}

\medskip\noindent\textbf{Partie IV --- \omterm{def:ml:targets-labels-and-sample-weights:meta}{Méta-étiquetage} et verdict.}
\begin{enumerate}[resume]
\item Que gagnent les transactions conservées à un seuil de 0.52, et que perd le portefeuille~?
\item Pourquoi le \omterm{def:ml:targets-labels-and-sample-weights:meta}{méta-étiquetage} relève-t-il du dimensionnement et non de l'alpha~?
\item Faut-il utiliser la probabilité pour dimensionner plutôt que pour filtrer~? Comment~?
\item Énoncer le \emph{résultat nommé}~: le nombre effectif d'observations face au nombre nominal, la $t$-statistique
avant et après, et le gain d'AUC de la pondération par l'unicité par rapport au bagging simple, avec sa dispersion.
\item Que deviendrait l'échantillon effectif avec des étiquettes à 20 jours~?
\item Comment un journal de recherche doit-il consigner la taille d'un \omterm{def:ml:why-financial-machine-learning-is-different:sets}{ensemble d'entraînement}~?
\item Que vient faire la purge (livre~7) avec l'unicité~?
\item En une phrase~: qu'est-ce qu'une étiquette~?
\end{enumerate}
\end{problem}

\section{Questions d'entretien}

\begin{interviewq}[$\star$ \iqroles{chercheur, ingénieur ML}]\label{iq:ml:targets-labels-and-sample-weights:1}
Qu'est-ce qui ne va pas dans le fait de labelliser chaque jour par le signe du rendement des 20 jours suivants~?
\end{interviewq}

\begin{interviewq}[$\star$ \iqroles{chercheur}]\label{iq:ml:targets-labels-and-sample-weights:2}
Pouvez-vous expliquer la méthode à triple barrière en deux minutes~?
\end{interviewq}

\begin{interviewq}[$\star\star$ \iqroles{chercheur, ingénieur ML}]\label{iq:ml:targets-labels-and-sample-weights:3}
Vous disposez d'un million de lignes d'étiquettes à 10 jours, échantillonnées chaque jour sur 400 actions pendant dix ans. Combien d'observations indépendantes
avez-vous, à peu près, et comment cela modifie-t-il votre choix de modèle~?
\end{interviewq}

\begin{interviewq}[$\star\star$ \iqroles{chercheur, trader}]\label{iq:ml:targets-labels-and-sample-weights:4}
Qu'est-ce que le \omterm{def:ml:targets-labels-and-sample-weights:meta}{méta-étiquetage}, et quand est-il utile~?
\end{interviewq}

\begin{interviewq}[$\star\star$ \iqroles{ingénieur ML}]\label{iq:ml:targets-labels-and-sample-weights:5}
Le score hors sac d'une forêt aléatoire sur des étiquettes qui se chevauchent est bien meilleur que son score sur une période de test ultérieure. Pourquoi~?
\end{interviewq}

\begin{interviewq}[$\star\star\star$ \iqroles{chercheur}]\label{iq:ml:targets-labels-and-sample-weights:6}
Pouvez-vous montrer que la moyenne de $n$ sommes de $h$ périodes qui se chevauchent, d'un bruit blanc, a environ $h$ fois la variance que vous calculeriez
si vous les traitiez comme indépendantes.
\end{interviewq}
